% %%%%%%%%%%%%%%%%%%%%%%%%%%%%%%%%%%%%%%%%%%%%%%%%%%%%%%%%%%%%%%%%%%%%%%%%%%%%%%%%%
%
%	cap#.tex 
%  -----------
%
%	Neste arquivo você escreve o capítulo. Deve sempre iniciar com o ambiente
%	\chapter{Nome do Capitulo}
%
% %%%%%%%%%%%%%%%%%%%%%%%%%%%%%%%%%%%%%%%%%%%%%%%%%%%%%%%%%%%%%%%%%%%%%%%%%%%%%%%%%
%
\chapter{Introdução}\label{introducao}
% ----------------------------------------------------------

Tema da pesquisa: Deformações induzidas pelo processo construtivo defasado de túneis gêmeos profundos com galerias transversais considerando o comportamento instantâneo e diferido no tempo.
  
% ---
\section{Objetivos}
% ---

O principal objetivo do presente trabalho consiste em programar/aplicar modelos constitutivos para o maciço e para o revestimento, a fim de analisar o comportamento instantâneo e diferido no tempo através do processo de escavação defasado de túneis profundos interligados por galerias.

\begin{alineas}
	
	\item elaborar um script APDL através do software ansys; 
	
	\item empregar a USERMAT;
	
	\item analisar as convergências;
	
	\item objetivo específico 4.
	
\end{alineas}
% ---
\section{Metodologia}
% ---
Um parágrafo apresentando a metodologia do trabalho
% ---
\section{Justificativa}
% ---
Justificar a escolha da pesquisa (?)
% ---
\section{Delimitações}
% ---
Adota-se as seguintes delimitações para este trabalho:

\begin{alineas}
	
	\item Os túneis em estudo são considerados profundos, ou seja, são desprezadas as cargas provenientes de estruturas da superfície e também recalques induzidos pela escavação de túneis.
	
	\item O maciço pode apresentar descontinuidades e possuir um comportamento muito complexo. Portanto, para reduzir essa complexidade e mantendo a efetividade das simulações, neste trabalho admite-se um meio contínuo, homogêneo e isotrópico.
	
	\item Considerando a complexidade do estado de tensões internas do maciço devido à anisotropia, heterogeneidade e descontinuidade dos materiais, adotou-se um estado de tensões geostático-hidrostático, onde as tensões verticais e horizontais são iguais.
	
	\item velocidade de escavação constante
	
	\item modelo único de revestimento e com espessura constante 
	
	\item escavação é de face inteira
	
	\item citar os modelos utilizados (elastico, elastoplastico, viscoplastico...)
	
	\item evolução das deformações quase-estática (não se considera excitações dinâmicas)
	
\end{alineas}
% ---
\section{Delineamento do trabalho}
% ---
Apresentar os capítulos do trabalho: O capítulo 1 trata de...